\chapter{Crises, Ethics, Law, and the Social Dimensions of Finance}

\section{How crises propagate}
Financial crises often combine high leverage, fragile funding, concentrated exposures, opaque claims, and a shock. A loss can trigger margin calls. Margin calls force sales. Sales reduce prices and collateral values. Lower collateral tightens funding, which forces more sales. The feedback loop can transform a local loss into a system-wide liquidity crisis.

This sequence is a mechanism, not a universal chronology. Crises also arise from currency mismatch, sovereign stress, cyber incidents, fraud, operational failure, or policy mistakes. The important analytical question is which balance-sheet and network links transmit the shock.

\section{Runs and fire sales}
A bank or fund can be solvent in a long-run asset-value sense but unable to meet immediate withdrawals. If users believe others will withdraw, the belief can become self-fulfilling. Deposit insurance, central-bank liquidity, gates, redemption terms, capital, and resolution regimes alter incentives but also create moral hazard and distributional choices.

\section{Ethical duties}
Finance allocates opportunities and losses. Ethical practice includes honest disclosure, suitability, fair dealing, privacy, avoidance of market abuse, care with vulnerable users, and recognition of conflicts. A technically legal product can still exploit misunderstanding or power asymmetry. Fiduciary and professional duties vary by role and jurisdiction; the text does not replace the applicable rulebook.

\section{Models and responsibility}
A model can make an ethical choice look neutral by hiding who bears the tail. A credit model may ration access; a risk limit may liquidate positions during a panic; an incentive contract may reward short-term gains while leaving others with long-term losses. Analysts should ask: who is affected, who can contest the decision, which errors are asymmetric, and which assumptions are invisible to the user?

\section{Distribution and inequality}
Asset ownership, access to credit, inflation exposure, pensions, and insurance shape life chances. A higher expected return is not equally valuable to every household when liquidity needs and downside tolerance differ. Finance can broaden access, but products can also transfer fees and tail risk to people least able to bear them. Social evaluation therefore requires incidence analysis, not only aggregate market value.

\section{Sustainability and externalities}
Climate, biodiversity, labour conditions, and transition policy can affect cash flows, discount rates, physical assets, and legal liabilities. A sustainability label is not a single risk measure. Materiality, time horizon, scenario assumptions, data quality, and engagement policy must be stated. Greenwashing is both an information problem and a governance problem.

\section{Crisis learning}
Post-crisis reforms should be evaluated by mechanisms and incentives. More capital can absorb losses; more liquidity can reduce runs; central clearing can reduce bilateral exposures; disclosure can improve information. Each intervention can have boundary effects. The relevant question is not whether a reform sounds protective but whether it changes the failure mechanism without creating a larger one.

\source{The Basel Committee's post-crisis framework, IOSCO market-infrastructure principles, and central-bank financial-stability reviews provide institutional context. Current legal obligations depend on the relevant jurisdiction and date.}

\begin{exerciseblock}
\begin{enumerate}
  \item Trace a four-step fire-sale feedback loop.
  \item Distinguish insolvency from illiquidity.
  \item Give two ethical questions that a credit-scoring model should raise beyond predictive accuracy.
  \item Why is a sustainability label not by itself evidence of lower financial risk?
\end{enumerate}
\end{exerciseblock}

